\subsection{Bad-Character Rule}
\begin{frame}{Bad-Character 직관}
window 마지막 text character \(c=T[s+m-1]\)를 pattern의 오른쪽 끝 이전 rightmost \(c\)와 맞춘다.
\[
shift[c]=m-1-\max\{j<m-1:P[j]=c\},
\]
그런 \(j\)가 없으면 \(m\).
\end{frame}
\begin{frame}{\texttt{TIGER} Shift}
\centering\begin{tikzpicture}
\node[smhash]at(0,1){\only<1|handout:0>{\(T:5-1-0=4,\ I:3,\ G:2\)}\only<2|handout:1>{\(E:1,\ R:5,\ other:5\)}};
\node[smnote]at(0,0){마지막 \(R\)은 preprocessing loop에서 제외};
\end{tikzpicture}
\end{frame}
\begin{frame}{Pattern에 없는 Character}
\[
P=\texttt{TIGER},\qquad \text{window last}=\texttt{b}.
\]
\centering\begin{tikzpicture}
\node[smhash]at(0,.8){\only<1|handout:0>{\(\texttt{b}\notin P[0..3]\)}\only<2|handout:1>{\(shift[\texttt{b}]=5\): window 전체 jump}};
\node[smnote]at(0,-.3){skipped region에는 시작 가능한 occurrence가 없다};
\end{tikzpicture}
\end{frame}
\begin{frame}{Pattern 안쪽 Character}
\[
P=\texttt{TIGER},\qquad \text{window last}=\texttt{I}.
\]
\[
shift[\texttt{I}]=5-1-1=3.
\]
pattern의 \(I\)를 현재 window 끝의 \(I\)에 맞추도록 3칸 이동한다.
\end{frame}
\begin{frame}{Full Boyer--Moore Preview}
full algorithm은 bad-character shift와 good-suffix shift를 모두 계산하고 안전한 더 큰 shift를 사용한다.
\scheckpoint{Lecture 9의 구현과 trace는 good-suffix table이 없는 Boyer--Moore--Horspool이다.}
\end{frame}
